\documentclass[10pt]{article}
\usepackage[T1]{fontenc}
\usepackage[utf8]{inputenc}
\usepackage[french]{babel}
\usepackage{amsmath,amssymb,mathrsfs,mathpazo}
\usepackage{geometry}
\geometry{a4paper,margin=2cm}
\usepackage{graphicx,array,tabularx,tabularray,colortbl,multirow,multicol,booktabs,enumitem,pifont,ulem,fancybox,diagbox,cancel}
\usepackage{pstricks,pstricks-add,pst-plot,pst-tree,pst-node,pst-text,pst-func,pst-math,pst-bspline,pst-3dplot,pst-eucl,pst-solides3d}
\let\pscalculate\undefined
\usepackage{pst-bezier}
\usepackage{tikz,pgfplots,tkz-tab}
\usetikzlibrary{arrows,arrows.meta,patterns,calc,positioning,trees,decorations.markings,shapes.geometric,shapes.misc,shapes.arrows,intersections,decorations.pathmorphing,decorations.pathreplacing,backgrounds,shadows,babel}
\pgfplotsset{compat=1.18}
\usepgfplotslibrary{fillbetween}
\usepackage[np]{numprint}
\usepackage{listings,makecell,pas-tableur,twoopt}
\let\barre\undefined
\usepackage[pstricks]{tablvar}
\definecolor{darkgreen}{RGB}{0,100,0}
\providecommand{\red}{\color{red}}
\providecommand{\blue}{\color{blue}}
\providecommand{\green}{\color{green}}
\providecommand{\black}{\color{black}}
\providecommand{\white}{\color{white}}
\providecommand{\gray}{\color{gray}}
\providecommand{\cyan}{\color{cyan}}
\providecommand{\magenta}{\color{magenta}}
\providecommand{\yellow}{\color{yellow}}
\providecommand{\violet}{\color{violet}}
\providecommand{\orange}{\color{orange}}
\providecommand{\euro}{\texteuro}
\providecommand{\EUR}{\texteuro}
\providecommand{\ieme}{\textsuperscript{e}}
\providecommand{\degres}{\ensuremath{^\circ}}
\providecommand{\degree}{\ensuremath{^\circ}}
\providecommand{\vv}[1]{\overrightarrow{#1}}
\providecommand{\Coord}[1]{#1}
\providecommand{\decosix}{}
\providecommand{\starredbullet}{\ensuremath{\bullet}}
\pagestyle{empty}
\setlength{\parindent}{0pt}
\usepackage{adjustbox}\DeclareUnicodeCharacter{2236}{:}\DeclareUnicodeCharacter{2190}{\ensuremath{\leftarrow}}
\providecommand{\R}{\mathbb{R}}
\providecommand{\N}{\mathbb{N}}
\providecommand{\Z}{\mathbb{Z}}
\providecommand{\Q}{\mathbb{Q}}
\providecommand{\D}{\mathbb{D}}
\providecommand{\C}{\mathbb{C}}
\providecommand{\e}{{\mathrm{e}}}
\providecommand{\vect}[1]{\overrightarrow{#1}}
\providecommand{\vectt}[1]{\overrightarrow{\mathrm{#1}}}
\providecommand{\vvt}[1]{\overrightarrow{\mathrm{#1}}}
\providecommand{\np}[2][]{{#2~#1}}
\providecommand{\vv}[1]{{\overrightarrow{#1}}}
\renewcommand{\d}{{\mathrm{d}}}
\providecommand{\degres}{{^\circ}}
\providecommand{\degre}{{^\circ}}
\providecommand{\blue}{\color{blue}}
\providecommand{\red}{\color{red}}
\providecommand{\enskip}{\quad}
\providecommand{\llbracket}{\lbrack\!\lbrack}
\providecommand{\rrbracket}{\rbrack\!\rbrack}
\providecommand{\euro}{{\text{€}}}
\providecommand{\no}{{\text{n°}}}
\providecommand{\interval}[2]{{[#1;#2]}}
\providecommand{\intervFO}[2]{{[#1;#2[}}
\providecommand{\dotfill}{\ldots}

\providecommand{\ds}{\displaystyle}
\providecommand{\textding}[1]{#1}
\providecommand{\Oij}{\(\mathrm{O};\vec{\imath},\vec{\jmath}\)}
\providecommand{\Oijk}{\(\mathrm{O};\vec{\imath},\vec{\jmath},\vec{k}\)}
\providecommand{\Ouv}{\(\mathrm{O};\vec{u},\vec{v}\)}
\providecommand{\intervFF}[2]{\ensuremath{[#1;#2]}}
\renewcommand{\euro}{\ifmmode\text{\texteuro}\else\texteuro\fi}
\DeclareUnicodeCharacter{00B0}{\ifmmode^{\circ}\else\textdegree\fi}
\DeclareUnicodeCharacter{20AC}{\ifmmode\text{\texteuro}\else\texteuro\fi}
% Énoncés : Éducation nationale. Transcription et corrigés : APMEP, Denis Vergès et contributeurs. Découpage et classement : CoopMaths, Sébastien Lozano, Éric Elter et contributeurs. Adaptation de la présentation : Maths Explorer.
% CC BY-SA : https://coopmaths.fr/www/about/
% https://coopmaths.fr/alea/static/bac/2021/tex/bac_2021_06_asie_2.tex
\begin{document}
\noindent\textit{APMEP / CoopMaths. CC BY-SA. Présentation adaptée par Maths Explorer.}\par\medskip
\medskip

On considère un cube ABCDEFGH d'arête 8 cm et de centre $\Omega$.

Les points P, Q et R sont définis par $\vect{\text{AP}} = \dfrac{3}{4}\vect{\text{AB}},\, \vect{\text{AQ}} = ~\dfrac{3}{4}\vect{\text{AE}}$ et $\vect{\text{FR}} = \dfrac{1}{4}\vect{\text{FG}}$.

On se place dans le repère orthonormé $\left(\text{A}~;\vect{\imath},~\vect{\jmath},~\vect{k}\right)$ avec : $\vect{\imath} = \dfrac{1}{8}\vect{\text{AB}},\, \vect{\jmath}= \dfrac{1}{8}\vect{\text{AD}}$ et 

$\vect{k} = \dfrac{1}{8}\vect{\text{AE}}$.

\begin{center}
\psset{unit=0.85cm,arrowsize=2pt 4}
\begin{pspicture}(11,12)

\pspolygon(0,1.7)(6.5,0)(6.5,7.2)(0,8.9)
\uput[d](0,1.7){B}\uput[d](6.5,0){C}\uput[u](6.5,7.2){G}\uput[u](0,8.9){F}
\psline(6.5,0)(10,2.6)(10,9.8)(6.5,7.2)
\uput[r](10,2.6){D}\uput[u](10,9.8){H}
\psline(10,9.8)(3.5,11.5)(0,8.9)
\uput[u](3.5,11.5){E}
\psline[linestyle=dashed](0,1.7)(3.5,4.3)(3.5,11.5)
\uput[d](3.5,4.3){A}\uput[r](5,5.73){$\Omega$}
\psline[linestyle=dashed](3.5,4.3)(10,2.6)
\psline[linestyle=dashed](3.5,4.3)(6.5,7.2)
\psline[linestyle=dashed](3.5,11.5)(6.5,0)
\psdots(0.88,2.35)(3.5,9.7)(5,5.73)(1.625,8.475)
\uput[u](0.88,2.35){P}\uput[l](3.5,9.7){Q}\uput[u](1.625,8.475){R}
\uput[u](3.0625,3.975){$\vect{\imath}$}\uput[u](4.3125,4.0875){$\vect{\jmath}$}
\uput[l](3.5,5.2){$\vect{k}$}
\psline{->}(3.5,4.3)(3.0625,3.975) \psline{->}(3.5,4.3)(4.3125,4.0875) \psline{->}(3.5,4.3)(3.5,5.2)
\end{pspicture}
\end{center}

\textbf{Partie I}

\medskip

\begin{enumerate}
\item Dans ce repère, on admet que les coordonnées du point R sont (8~;~2~;~8). 

Donner les coordonnées des points P et Q.
\item Montrer que le vecteur $\vect{n}(1~;~-5~;~1)$ est un vecteur normal au plan (PQR).
\item Justifier qu'une équation cartésienne du plan (PQR) est $x - 5y + z - 6 = 0$.
\end{enumerate}

\bigskip

\textbf{Partie II}

\medskip

On note L le projeté orthogonal du point $\Omega$ sur le plan (PQR).

\medskip

\begin{enumerate}
\item Justifier que les coordonnées du point $\Omega$ sont (4~;~4~;~4).
\item Donner une représentation paramétrique de la droite $d$ perpendiculaire au plan (PQR) et passant par $\Omega$.
\item Montrer que les coordonnées du point L sont $\left(\dfrac{14}{3}~;~ \dfrac{2}{3}~;~\dfrac{14}{3}\right)$
\item Calculer la distance du point $\Omega$ au plan (PQR).
\end{enumerate}


\bigskip

\end{document}
